\documentclass [a4paper, 10pt]{scrartcl}

% Sprache + Encoding
\usepackage{ucs} % unicode
\usepackage[utf8x]{inputenc}
\usepackage{ngerman} % deutsche Sprache
\usepackage{ragged2e}

\parindent0pt
\parskip5mm plus1mm minus1mm

\setlength{\textwidth}{140mm}
\setlength{\textheight}{265mm}
\setlength{\oddsidemargin}{12mm}
\setlength{\evensidemargin}{12mm}
\setlength{\topmargin}{-20mm}
\setlength{\headheight}{0mm}
\setlength{\columnsep}{5mm}
\setlength{\footskip}{18mm}
\setlength{\headsep}{12mm}

\usepackage{german}
\usepackage{graphicx}
\usepackage{tabularx}
\usepackage{amsmath}
\usepackage{mdwlist}
\usepackage{tabularx}



\newcommand{\Referent}[2][]{
  \newcommand{\betreuertitel}{#1\ }
  \newcommand{\betreuername}{#2}        
}


\newcommand{\Mitarbeiter}[2][]{%
  \newcommand{\mitarbeitertitel}{#1\ }%
  \newcommand{\mitarbeitername}{#2}%
}

\newcommand{\Student}[2][]{%
  \newcommand{\studenttitel}{#1\ }%
  \newcommand{\studentname}{#2}%
}
\def\Matrikelnummer #1 {%
  \def\matrikelnr {#1}%
}

\def\Titel #1 {%
  \def\arbeittitel {#1}%
}

\def\Arbeitsart #1 {% 
  \def\art {#1}
}

\newenvironment{aufgaben}[0]{\textbf{Aufgaben:}\\\vspace*{-10mm}\begin{itemize*}}{\end{itemize*}}

\def\Termine #1#2#3 {%
  \def\beginn {#1} \def\zwischen {#2} \def\ende {#3}%
}

\def\Anfang{%
  \centerline{\resizebox{\linewidth}{!}{\includegraphics{isasheader.pdf}}}%
  % 
  \sffamily 
  \vspace*{5mm}%
  \centerline{\Large{%
      \textbf{\art:}}}\vspace*{-5mm}
  \begin{center}
    \Large{%
      \textbf{\arbeittitel}}
  \end{center}
  \bigskip
}			

\def\Aufgaben #1{%
  \def\aufgabenbesch {#1}
}

\def\Ende{%
  \vfill
  \hspace{-\parindent}%
  \begin{tabularx}{\textwidth}{lXlX}
    % \bf Titel: &\bf Effiziente Prädiktion nichtlinearer Systeme basierend auf der Approximation von Transitionsdichten mittels Gaussian Mixtures\\
    \textbf{Bearbeiter:} &\textbf{\studenttitel\studentname} &
    \textbf{Matrikelnummer: \matrikelnr}
    \\\ \\
    \textbf{Betr. Mitarbeiter:}&\multicolumn{2}{l}{\textbf{\mitarbeitertitel\mitarbeitername}}\\
    \textbf{Referent:}& \multicolumn{2}{l}{\textbf{\betreuertitel\betreuername}}\\
    \ \\\end{tabularx}
  \begin{tabular}{lr}
    \textbf{Beginn:} &\textbf{\beginn}\\
    \textbf{Zwischenvortrag:} & \textbf{\zwischen}\\
    \textbf{Abgabe:} & \textbf{\ende}
  \end{tabular}
  \par
  \vfill


  Karlsruhe, den \today\par
  \vspace{1.5cm}
  \noindent\begin{tabularx}{1.0\linewidth}[h]{cXcXc}
    \underline{\hspace{4cm}}&\ &\underline{\hspace{4.5cm}}
    &\ &\underline{\hspace{4cm}}\\
    \emph{\studentname} && \emph{\mitarbeitername} && \emph{\betreuername}
  \end{tabularx}  

}

\Arbeitsart{Masterarbeit}
\Titel{Bessere Rekonstruktion bei weniger Messungen durch Kombination von Kalman-Filter und Compressive-Sensing}
\Ziel{Die bessere Signalrekonstruktion bei weniger Messungen durch Kombination von Kalman-Filter und Compressive-Sensing}

\Student[B.Sc.]{Zhao, Haibin}
\Matrikelnummer{2247796}
\Referent[Prof. Dr.-Ing.]{Uwe D. Hanebeck}
\Mitarbeiter[Dr.-Ing.]{Noack, Benjamin}
\MitarbeiterB[M.Sc.]{Funk, Christopher}

% Beginn Arbeit / Zwischenvortrag / Abgabe
\Termine{15. Mai 2020}{$\approx$15. August 2020}{15. November 2020}

% Datum für Unterschrift:
%\renewcommand{\today}{20. Mai 2007}

% Ganttchart
\Ganttchart{gantt.png}

\begin{document}
\Anfang

Wenn ein Signal dünnbesetzt ist, kann man Compressive-Sensing (CS) anwenden, um das Signal durch weniger Messungen zu rekonstruieren, als nach dem Nyquist-Shannon Theorem gefordert sind. CS ist jedoch nicht für nicht-dünnbesetzte Signale geeignet. Daher muss man ein neues Verfahren aufstellen, um CS anzuwenden. Mögliche Maßnahmen sind:
\begin{enumerate}
\item Betrachten der Änderung des Signals, anstatt des Signals selbst. Wenn die Änderung des Signals dünnbesetzt ist, kann man CS anwenden, um die Änderung der Koeffizienten zu schätzen. Im Vergleich zur Forderung dünnbesetzter Signale ist dies eine schwächere Forderung. Insbesondere bei kleinem zeitlichem Abtastintervall, kann sie häufig erfüllt werden.

\item Reduzieren der Dimension des zu rekonstruierenden Signals. Eine bestimmte Anzahl von Sensoren werden eingesetzt, für jeden Zeitschritt werden einige Sensoren zufällig aktivert. Auf diese Weise wird CS nicht auf das ganze Signal, sondern auf alle möglichen Sensorstellen angewendet. Dadurch kann die Dimension deutlich verringert werden, wodurch die Dünnbesetztheit des Signals erhöht wird.

\item Kombinieren von CS mit einem Systemmodell (Kalman-Filter). Das rekonstruierte Signal aus CS kann als Pseudo-Messungen in einem Kalman-Filter eingesetzt werden. Gleichzeitig werden die Messunsicherheiten dynamisch bestimmt: bei den realen Messungen  (Messungen aus Sensoren) werden niedrige Messunsicherheiten und bei den Pseudo-Messungen (rekonstruierten Werten aus CS) hohe Messunsicherheiten eingestellt.

\item Sparse Coding Representation (SCR). Mithilfe des maschinellen Lernens werden Basisfunktionen, mit denen das Signal dünnbesetzt dargestellt kann, gesucht.
\end{enumerate}

Da im obigen Prozess nur die Änderung des Signals geschätzt wird, führt dies zu einer kontinuierlichen Akkumulation von Fehlern. Daher sind einige Methoden erforderlich, um den akkumulierten Fehler zu reduzieren, z.B.
\begin{enumerate}
\item Die Basiskoeffizienten des vom Kalman-Filter geschätzten Signals werden im nächsten Zeitschritt für das CS-Verfahren verwendet.

\item Die Prädiktion des Kalman-Filters wird zur Berechnung der Änderung der Koeffizienten genutzt.
\end{enumerate}

Darüber hinaus kann man das CS-Verfahren in jedem Zeitschritt iterativ durchführen. D.h. man berechnet die Änderung der Koeffizienten iterativ, bis eine Endbedingung erreicht ist. 
\\
\\
\\
\\
\\
\\
\\
\\
\\

\begin{aufgaben}
\item Literaturüberblick über verwandte Ansätze zur Zustandsschätzung geben.
\item Erstellen von Testumgebung (Temperaturverteilung)
  \begin{itemize}
  \item Modalanalyse
  \item Finite Differenz Methode
  \end{itemize}
\item Grundlagen für Kalman-Filter und Compressive Sensing
\item Kalman-Filtered Compressive Sensing
  \begin{itemize}
  \item zufällige Messungen + Kalman-Filter
  \item dünnbesetzte Änderung
  \item Dimensionsreduktion
  \item Sparse Coding Representation
  \item Integration von KF und CS
  \item Iteratives Verfahren
  \end{itemize}
\item Abschluss
  \begin{itemize}
  \item Präsentation
  \item Masterarbeit schreiben
  \end{itemize}
\end{aufgaben}%
\Ende
\end{document}
